% ===================== 附录 =====================
\clearpage
\SetStage{mainNavy}{附录}{曲库总表}

\pdfbookmark[1]{附录 · 曲库总表}{appA}
\pagehead{附录 · 曲库总表}

大纲中出现的全部曲目及其所在课次如下，方便按曲目倒查学习进度。

\vspace{4mm}
\noindent\begin{tabularx}{\textwidth}{@{}c l l X@{}}
\toprule
& {\sffamily\bfseries 课次} & {\sffamily\bfseries 出现方式} & {\sffamily\bfseries 曲目} \\
\midrule
\stagesq{stageA} & 第 9 课  & 选曲     & 《Fade to Black》 \\
\stagesq{stageA} & 第 10 课 & 曲目练习 & 《不再犹豫》前奏扫弦 \\
\stagesq{stageB} & 第 12 课 & 选曲     & 《Fade to Black》 \\
\stagesq{stageB} & 第 13 课 & 曲目练习 & 谢天笑《循环的太阳》 \\
\stagesq{stageB} & 第 15 课 & 曲目练习 & 《天鹅之死》前奏 \\
\stagesq{stageB} & 第 17 课 & 选曲     & 《Fade to Black》副歌，或《Du Hast》前奏 \\
\stagesq{stageB} & 第 18 课 & 曲目练习 & 谢天笑《把夜晚染黑》前奏分解与 Solo \\
\stagesq{stageB} & 第 20 课 & 结课曲目 & 《One》 \\
\stagesq{stageC} & 第 24 课 & 曲目练习 & 乔伊的《Open Fire》 \\
\stagesq{stageC} & 第 25 课 & 曲目练习 & 乔伊的《Open Fire》 \\
\stagesq{stageC} & 第 27 课 & 长期作业 & 《Fade to Black》前奏 Solo \\
\stagesq{stageD} & 第 31--35 课 & 课程标题 & 《加州旅馆》Solo 学习（一）至（五） \\
\bottomrule
\end{tabularx}

\vspace{2mm}
\noindent{\footnotesize\color{inkGray!75}行首色块对应课程所属阶段（图例见「本册使用说明」）。
全书图示为按大纲内容绘制的原创矢量图及同风格 AI 生成插画（手型与舞台示范），无外部图片授权事项。}
